\hspace{3mm}

\centerline{\bf \Huge Тренировочная Олимпиада}

\hspace{3mm}

\problem{1} На доске написано число 3 728 954 106. Зачеркните в нем три цифры так, чтобы оставшиеся цифры в том же порядке образовали как можно меньшее число.

\hspace{3mm}

\problem{2} Андрей отметил в календаре 6 последовательных дней июня и за- помнил, на какие числа месяца пришлись эти дни. Потом он выписал эти 6 последовательных чисел в порядке убывания. Павлик перемно- жил первые четыре из этих чисел, Андрей — все, кроме двух крайних, а Майя — последние четыре. У Павлика и Андрея последние цифры результатов совпали, а у Майи получилась другая последняя цифра. Какая? Найдите все возможные ответы и докажите, что других нет.

\hspace{3mm}

\problem{3} В записи ***5 : 11 = ** замените звездочки цифрами так, чтобы получилось верное равенство.

\hspace{3mm}

\problem{4} В школе работают 5 кружков, в них занимаются 7, 10, 13, 15, 25 школьников (каждый школьник посещает только один кружок). Извест- но, что у каждого школьника в этих кружках есть не менее двух одно- классников. Докажите, что в хотя бы в одном кружке есть два школь- ника из одного класса.

\hspace{3mm}

\problem{5} За большим круглым столом сидят 200 человек. Каждый из них либо рыцарь, либо лжец, либо чудак. Рыцарь всегда говорит правду, лжец все- гда лжет. Чудак говорит правду, если справа от него сидит лжец; ложь, если справа от него сидит рыцарь; все что угодно, если справа от него чудак. Каждый сказал: «Слева от меня сидит лжец». Сколько всего лже- цов? Перечислите все возможные ответы и докажите, что других нет. 